% Procedencia: RESULTADO_RECUPERADO del capítulo 24 de la integral de 2249
% páginas y de FORMA_ELIPTICA_RADIO_MODULAR_ACCION_ORIENTADA.md (D04).
% Delta expositivo: dominio del operador diagonal escrito completamente y
% recuperación del entero en las dos cartas cocientadas hecha explícita.
% Dependencia externa de etiqueta: sec:ellipse. Bibliografía: elipse,
% hmtedicionintegral, flm, borcherds. No requiere macros nuevas.

\section{De la elipse de acción a la dualidad T}
\label{sec:ii-dualidad-t}

La reciprocidad de la elipse proporciona una coordenada de forma para el
intercambio entre valor visible y memoria de frontera. El antecedente es el
mismo estado enriquecido de APP--TRIT--TPK: la APP conserva las hojas de
suma y producto con sus residuos y cocientes; el TRIT conserva su orientación;
el TPK transporta hoja, acarreo, ruta y memoria. Las salidas angulares y de
acción ya construidas sobre la estructura discreta conjunta del continuo
determinan después la elipse. Esta sección compone sus lectores posteriores:
forma elíptica, intercambio reticular, realización de momento--enrollamiento
y evaluación modular \cite{elipse,hmtedicionintegral}.

El punto de corte se sitúa, por tanto, después de la generación de
\(A,C^*,\hbar\). Ningún radio prescrito ni constante de cuerda selecciona
esas salidas. La memoria de frontera se conserva primero en el par
residuo--cociente y se transporta luego mediante los mapas que se demuestran
a continuación.

\subsection{La forma de acción determina el radio normalizado}

Retomamos la coordenatización de la sección~\ref{sec:ellipse}. Sean \(A\ne0\),
\(\hbar>0\) y \(|C^*/A|<1\), con ambos ángulos expresados en la misma
unidad. Los semiejes permanecen etiquetados:
\begin{equation}
 \eta_{\rm el}=\operatorname{artanh}(C^*/A),\qquad
 a_S=\sqrt{2\hbar}\,e^{\eta_{\rm el}/2},\qquad
 b_S=\sqrt{2\hbar}\,e^{-\eta_{\rm el}/2}.
 \label{eq:ii-t-elipse}
\end{equation}
Las coordenadas equilibradas de fase y sus semiejes tienen unidad de raíz
de acción. El radio que sigue es adimensional.

\begin{proposition}[Radio de forma y reciprocidad etiquetada]
\label{prop:ii-t-radio}
La coordenatización rectangular de la elipse determina de manera única
\begin{equation}
 R_0=\sqrt{\frac{b_S}{a_S}}
 =e^{-\eta_{\rm el}/2}
 =\left(\frac{A-C^*}{A+C^*}\right)^{1/4}>0,
 \qquad \tau_{\rm el}=iR_0^2.
 \label{eq:ii-t-radio}
\end{equation}
Conserva la razón angular firmada mediante
\begin{equation}
 \frac{C^*}{A}=\frac{1-R_0^4}{1+R_0^4}.
 \label{eq:ii-t-recuperacion-angular}
\end{equation}
El intercambio de los semiejes etiquetados transforma
\((\eta_{\rm el},R_0,\tau_{\rm el})\) en
\((-\eta_{\rm el},R_0^{-1},-1/\tau_{\rm el})\), dejando fija
la escala \(a_Sb_S=2\hbar\).
\end{proposition}

\begin{proof}
La división de los semiejes de~\eqref{eq:ii-t-elipse} da
\(b_S/a_S=e^{-\eta_{\rm el}}\). La aplicación \(R\mapsto R^2\)
es biyectiva en \(\mathbb R_{>0}\), de modo que su raíz positiva es
única. Para \(x=C^*/A\in(-1,1)\),
\(2\operatorname{artanh}x=\log((1+x)/(1-x))\); por ello
\(R_0^4=(1-x)/(1+x)\). Despejar \(x\) prueba
\eqref{eq:ii-t-recuperacion-angular}. Intercambiar \(a_S,b_S\)
invierte su razón y cambia el signo de \(\eta_{\rm el}\). Finalmente,
\(-1/(iR_0^2)=iR_0^{-2}\), mientras el producto de los semiejes
permanece igual a \(2\hbar\).
\end{proof}

La razón angular se recupera; los dos ángulos absolutos y su historia
generadora siguen residiendo en el estado enriquecido. El punto \(R_0=1\)
corresponde a la forma circular \(C^*/A=0\), no a la identificación de
todas las orientaciones o de todas las historias.

\subsection{Intercambio reticular y equivalencia de operadores}

\begin{definition}[Dominio estable de la lectura dual]
\label{def:ii-t-reticular}
Sobre \(\mathscr D_0=\mathbb Z^2\times\mathbb R_{>0}\) definimos
\begin{equation}
 E_{R_0}(m,w)=\frac{m^2}{R_0^2}+w^2R_0^2,\qquad
 \mathsf T_0(m,w,R_0)=(w,m,R_0^{-1}).
 \label{eq:ii-t-energia}
\end{equation}
La primera coordenada prolonga la lectura visible y la segunda la de
cociente. El dominio entero completo permite intercambiarlas sin exigir
que el cociente pertenezca a una sección de nueve representantes.
\end{definition}

Para el radio de~\eqref{eq:ii-t-radio}, la matriz de esta forma es
\begin{equation}
 D(\eta_{\rm el})=
 \begin{pmatrix}e^{\eta_{\rm el}}&0\\0&e^{-\eta_{\rm el}}\end{pmatrix}
 =\frac{1}{2\hbar}
 \begin{pmatrix}a_S^2&0\\0&b_S^2\end{pmatrix}.
 \label{eq:ii-t-matriz}
\end{equation}
En efecto, \(R_0^{-2}=e^{\eta_{\rm el}}\) y
\(R_0^2=e^{-\eta_{\rm el}}\). La ecuación de la elipse
\(Q^2/a_S^2+P^2/b_S^2=1\), multiplicada por \(2\hbar\), utiliza
la matriz inversa \(D(-\eta_{\rm el})\). Así quedan diferenciadas
la matriz de semiejes cuadrados y la matriz de la ecuación elíptica.

\begin{theorem}[Dualidad reticular y conjugación unitaria]
\label{thm:ii-t-unitaria}
La transformación \(\mathsf T_0\) es involutiva y satisface
\begin{equation}
 E_{R_0^{-1}}(w,m)=E_{R_0}(m,w).
 \label{eq:ii-t-invariancia}
\end{equation}
En \(\mathcal H_{\rm lat}=\ell^2(\mathbb Z^2)\), sea \(H(R_0)\)
el operador diagonal de valores \(E_{R_0}(m,w)\), con dominio
\[
 \mathcal D(H(R_0))=
 \left\{c\in\ell^2(\mathbb Z^2):
 \sum_{m,w\in\mathbb Z}E_{R_0}(m,w)^2|c_{m,w}|^2<\infty\right\}.
\]
Entonces \(H(R_0)\) es autoadjunto y no negativo. El intercambio
\(U_T|m,w\rangle=|w,m\rangle\) es una involución unitaria que
transporta los dominios y satisface
\begin{equation}
 U_T H(R_0)U_T^{-1}=H(R_0^{-1}).
 \label{eq:ii-t-conjugacion}
\end{equation}
\end{theorem}

\begin{proof}
Intercambiar dos veces los enteros e invertir dos veces el radio devuelve
el punto inicial. Además,
\[
 E_{R_0^{-1}}(w,m)=w^2R_0^2+m^2R_0^{-2}=E_{R_0}(m,w).
\]
El dominio diagonal contiene las sucesiones de soporte finito y es denso.
La multiplicación por una sucesión real es simétrica en ese dominio.
Si \(c\) pertenece al dominio del adjunto, la evaluación contra cada
vector de la base obliga a que el adjunto tenga coordenadas
\(E_{R_0}(m,w)c_{m,w}\). Su pertenencia a \(\ell^2\) equivale
exactamente a la condición que define \(\mathcal D(H(R_0))\).
Luego el operador y su adjunto coinciden. Sus valores diagonales no
negativos prueban la positividad.

El intercambio permuta biyectivamente la base ortonormal, por lo que es
unitario y su inverso es él mismo. Para toda sucesión \(c\),
\[
 \sum_{m,w}E_{R_0^{-1}}(m,w)^2|(U_Tc)_{m,w}|^2
 =\sum_{m,w}E_{R_0}(m,w)^2|c_{m,w}|^2.
\]
Esta igualdad prueba el transporte exacto de dominios. Finalmente,
\((U_TH(R_0)U_T^{-1}c)_{m,w}=E_{R_0}(w,m)c_{m,w}
=E_{R_0^{-1}}(m,w)c_{m,w}\); así queda demostrada
\eqref{eq:ii-t-conjugacion} en todo el dominio, no sólo sobre la base.
\end{proof}

La orientación añade información que la energía no distingue. En el plano
de fase, el intercambio \(S_2(Q,P)=(P,Q)\) y el giro
\(J_2(Q,P)=(-P,Q)\) transportan la misma forma cuadrática, pero
\(S_2^*(dQ\wedge dP)=-dQ\wedge dP\) y
\(J_2^*(dQ\wedge dP)=dQ\wedge dP\). Con
\(\lambda=P\,dQ\), se obtienen directamente
\[
 S_2^*\lambda=d(PQ)-\lambda,\qquad
 J_2^*\lambda=\lambda-d(PQ).
\]
La integración sobre una curva cerrada da acciones opuestas bajo \(S_2\)
e iguales bajo \(J_2\). Por tanto, el intercambio de índices de
\eqref{eq:ii-t-energia} y el transporte simpléctico de una trayectoria
son mapas distintos, aun cuando publiquen la misma energía.

\subsection{Secciones nonádicas y transporte del acarreo}

Para \(n\in\mathbb Z\), la sección positiva escribe
\[
 r_9^+(n)=1+((n-1)\bmod9),\qquad
 Q_9^+(n)=\frac{n-r_9^+(n)}9,\qquad n=r_9^+(n)+9Q_9^+(n).
\]
El resto de \(n-1\) se toma en \(\{0,\ldots,8\}\). El cambio a
la sección de restos \(0,\ldots,8\) es
\begin{equation}
 C(r,Q)=\bigl(r-9\varepsilon,Q+\varepsilon\bigr),\qquad
 \varepsilon=\begin{cases}1,&r=9,\\0,&r\ne9.\end{cases}
 \label{eq:ii-t-acarreo}
\end{equation}
Así, \(r+9Q\) permanece fijo. Para \(n=9\), las dos coordenadas
son \((9,0)\) y \((0,1)\): sus energías no corregidas son
\(81/R_0^2\) y \(R_0^2\). La igualdad del entero no implica
igualdad de esas inserciones energéticas. La forma covariante siguiente
resuelve el cambio de sección sin suprimir el acarreo.

\begin{definition}[Dos coordenatizaciones cocientadas de la lectura nonádica]
\label{def:ii-t-cociente}
Sean \(L_9=\mathbb Z(9,-1)\), \(S(m,w)=(w,m)\) y
\(L_9^\vee=SL_9=\mathbb Z(-1,9)\). Para
\(L\in\{L_9,L_9^\vee\}\) definimos
\[
 \mathcal E_{R_0}^{L}([x]_L)
 =\min_{\ell\in L}E_{R_0}(x+\ell).
\]
El transporte entre coordenatizaciones es
\[
 \mathsf T_{\rm cov}:
 (\mathbb Z^2/L_9)\times\mathbb R_{>0}
 \longrightarrow
 (\mathbb Z^2/L_9^\vee)\times\mathbb R_{>0},\qquad
 ([x]_{L_9},R_0)\longmapsto([Sx]_{L_9^\vee},R_0^{-1}).
\]
\end{definition}

\begin{theorem}[Covariancia con memoria de cruce]
\label{thm:ii-t-cociente}
Las energías y el transporte anteriores están bien definidos. El cambio
de sección~\eqref{eq:ii-t-acarreo} conserva la clase en
\(\mathbb Z^2/L_9\), y
\begin{equation}
 \mathcal E_{R_0^{-1}}^{L_9^\vee}([Sx]_{L_9^\vee})
 =\mathcal E_{R_0}^{L_9}([x]_{L_9}).
 \label{eq:ii-t-covariancia}
\end{equation}
El intercambio de vuelta recupera la clase y el radio iniciales. Las
aplicaciones de coordenadas enteras \(N([x])=x_1+9x_2\) y
\(N^\vee([y])=9y_1+y_2\) son isomorfismos con \(\mathbb Z\)
y satisfacen \(N^\vee([Sx])=N([x])\).
\end{theorem}

\begin{proof}
Sustituir \(x\) por \(x+\ell_0\), con \(\ell_0\in L\),
permuta los candidatos al mínimo. Éste existe: para \(L=L_9\),
\[
 E_{R_0}(x+k(9,-1))
 =\left(\frac{81}{R_0^2}+R_0^2\right)k^2
  +2\left(\frac{9x_1}{R_0^2}-x_2R_0^2\right)k+E_{R_0}(x),
\]
que tiende a infinito con \(|k|\); para \(L_9^\vee\), el
coeficiente principal es \(R_0^{-2}+81R_0^2>0\). El cambio
\(C(x)-x=-\varepsilon(9,-1)\) pertenece a \(L_9\).
Si \(x'-x\in L_9\), entonces \(Sx'-Sx\in L_9^\vee\),
lo que prueba que el transporte está bien definido. Por
\eqref{eq:ii-t-invariancia},
\[
 \min_{\ell\in L_9}E_{R_0^{-1}}(Sx+S\ell)
 =\min_{\ell\in L_9}E_{R_0}(x+\ell).
\]
Esto demuestra~\eqref{eq:ii-t-covariancia}; \(S^2=I\) prueba
el retorno entre las dos coordenatizaciones.

Por último, el homomorfismo \(x\mapsto x_1+9x_2\) es
sobreyectivo y su núcleo es exactamente \(L_9\), pues
\(x_1=-9x_2\). Análogamente, el núcleo de
\(y\mapsto9y_1+y_2\) es \(L_9^\vee\). Sus aplicaciones
inducidas en los cocientes son, por tanto, isomorfismos. Sustituir
\(y=Sx\) prueba la identidad final.
\end{proof}

El entero recuperado determina sus residuos y cocientes una vez fijada
la sección. El representante que minimiza la energía puede no ser ese
representante canónico. El cociente es una representación covariante de la
lectura nonádica; el estado fuente conserva además hoja, ruta, orientación,
frontera, memoria y supervivencia. Tampoco se reemplaza \(L_9^\vee\)
por \(L_9\): el intercambio es entre dos coordenatizaciones transportadas.

\subsection{Momento, enrollamiento y escala de la realización}

La realización circular introduce una coordenatización dimensional explícita.
Denotemos por \(R>0\) su radio y por \(\alpha'>0\) un parámetro
de dimensión longitud al cuadrado. Este símbolo es distinto de la salida
adimensional \(\alpha_{\rm HMT}\). Fijar la coordenatización no altera el
generador angular ni el radio de forma. La normalización y su inversa son
\[
 \mathcal N_{\alpha'}(m,w,R)
 =\left(m,w,\frac{R}{\sqrt{\alpha'}}\right),\qquad
 R=\sqrt{\alpha'}\,R_0.
\]
Los índices se leen ahora como número de momento y número de
enrollamiento; las componentes de modo cero, expresadas en unidades de
inversa de longitud, son
\begin{equation}
 p_L=\frac mR+\frac{wR}{\alpha'},\qquad
 p_R=\frac mR-\frac{wR}{\alpha'},\qquad
 E_{\rm str}=\frac{p_L^2+p_R^2}{2}
 =\frac{m^2}{R^2}+\frac{w^2R^2}{\alpha'^2}.
 \label{eq:ii-t-quirales}
\end{equation}

\begin{theorem}[Conjugación dimensional de la dualidad T]
\label{thm:ii-t-cuerda}
La transformación
\[
 \mathsf T_{\alpha'}(m,w,R)=\left(w,m,\frac{\alpha'}R\right)
\]
es una involución y satisface
\begin{equation}
 \mathcal N_{\alpha'}\mathsf T_{\alpha'}
 =\mathsf T_0\mathcal N_{\alpha'},\qquad
 E_{\rm str}=\alpha'^{-1}E_0\circ\mathcal N_{\alpha'},
 \label{eq:ii-t-normalizacion}
\end{equation}
donde \(E_0(m,w,R_0)=E_{R_0}(m,w)\). Conserva \(E_{\rm str}\),
deja fijo \(p_L\) e invierte el signo de \(p_R\). Para datos
oscilatorios \(N_L,N_R\) e interceptos \(a_L,a_R\) fijos,
conserva también la condición
\begin{equation}
 \frac{\alpha'}4(p_L^2-p_R^2)+N_L-N_R-a_L+a_R=0,
 \quad\text{es decir}\quad
 N_L-N_R=a_L-a_R-mw.
 \label{eq:ii-t-nivel}
\end{equation}
\end{theorem}

\begin{proof}
Dos aplicaciones intercambian dos veces los índices y envían
\(R\) a \(\alpha'/(\alpha'/R)=R\). Asimismo,
\[
 \mathcal N_{\alpha'}\mathsf T_{\alpha'}(m,w,R)
 =\left(w,m,\frac{\sqrt{\alpha'}}R\right)
 =\mathsf T_0\mathcal N_{\alpha'}(m,w,R).
\]
La sustitución en la forma normalizada da
\[
 E_0\left(m,w,\frac R{\sqrt{\alpha'}}\right)
 =\frac{\alpha'm^2}{R^2}+\frac{w^2R^2}{\alpha'}
 =\alpha'E_{\rm str}(m,w,R).
\]
Para \(R'=\alpha'/R\), los modos transformados son
\[
 p'_L=\frac{wR}{\alpha'}+\frac mR=p_L,\qquad
 p'_R=\frac{wR}{\alpha'}-\frac mR=-p_R.
\]
Por ello se conservan tanto su suma de cuadrados como su diferencia.
La expansión directa da \(p_L^2-p_R^2=4mw/\alpha'\), que
transforma la primera ecuación de~\eqref{eq:ii-t-nivel} en la segunda.
El intercambio conserva \(mw\); con los otros datos fijos, conserva
la condición completa.
\end{proof}

La fórmula abreviada \(N_L-N_R=nw\) exige, con estas orientaciones,
\(a_L=a_R\) y \(n=-m\). Mantener el signo de
\eqref{eq:ii-t-nivel} conserva el diccionario entre las dos lecturas.
El resultado demostrado comprende la conjugación de radios, los modos
cero, la energía y la compatibilidad de niveles declarada. La interpretación
dimensional identifica las coordenadas mediante su coordenatización; no convierte
\(R_0\) por sí solo en una longitud ni deduce una teoría de campos
completa de una igualdad de energías.

\subsection{La evaluación modular y el carácter de Moonshine}

El corredor excepcional ya desarrollado conserva la incidencia del mismo
antecedente hasta \(V^\natural\). Su carácter graduado tiene la expresión
\(J(\tau)=j(\tau)-744\), con la invariancia modular empleada en
Moonshine \cite{flm,borcherds}. Esta \(J\) es una función en el
semiplano superior \(\mathbb H\); no es el giro \(J_2\), el
intercambio \(S\) ni la holonomía nonádica \(\Gamma_9\).

\begin{proposition}[Semiconjugación modular del radio]
\label{prop:ii-t-modular}
Sean \(\iota:\mathbb R_{>0}\to\mathbb H\),
\(\iota(R_0)=iR_0^2\), y
\(S_{\rm mod}(\tau)=-1/\tau\). Entonces
\begin{equation}
 S_{\rm mod}\circ\iota(R_0)=\iota(R_0^{-1}),\qquad
 J(iR_0^2)=J(iR_0^{-2}).
 \label{eq:ii-t-modular}
\end{equation}
En particular, las dos formas etiquetadas de la elipse tienen igual
valor de \(J\).
\end{proposition}

\begin{proof}
La primera identidad es
\(-1/(iR_0^2)=iR_0^{-2}\). La matriz
\(\left(\begin{smallmatrix}0&-1\\1&0\end{smallmatrix}\right)\)
pertenece a \(\mathrm{SL}_2(\mathbb Z)\) y actúa en
\(\mathbb H\) como \(S_{\rm mod}\). La invariancia modular de
\(J\), establecida en su construcción excepcional, da
\(J(S_{\rm mod}\tau)=J(\tau)\). Aplicándola a
\(\tau=\iota(R_0)\) se obtiene la segunda identidad.
La proposición~\ref{prop:ii-t-radio} identifica esos dos argumentos
con los módulos de las elipses de formas opuestas.
\end{proof}

La composición conserva así una cadena explícita:
\[
 (A,C^*,\hbar)\longmapsto(a_S,b_S)
 \longmapsto R_0\longmapsto iR_0^2\longmapsto J(iR_0^2).
\]
La inversión de forma se transporta a inversión de radio y a
\(S_{\rm mod}\); la acción orientada y el cambio de sección conservan
sus mapas propios. La función \(J\) asigna un valor a una clase modular, mientras
el marcado de los ejes distingue sus dos presentaciones y el estado TPK
conserva su genealogía. Ésta es la relación demostrada entre la elipse de
acción, la dualidad T y Moonshine: una composición con dominios e
identidades de transporte, cuyo valor modular no sustituye ni el operador
reticular ni la construcción de \(V^\natural\).
